% -*- coding: utf-8 -*-
% !TEX program = xelatex
% !TEX encoding = utf-8

\documentclass{nkuthesis}

% 参考文献
\usepackage[maxnames=3,minnames=3,sorting=none,style=nkthesis]{biblatex}
\addbibresource{reference.bib}

% 数学环境
\usepackage{amssymb}
\usepackage{amsmath}
\usepackage{newtxmath}
\newtheorem{Theorem}{\hspace{2em}定理}[section]
\newtheorem{Lemma}[Theorem]{\hspace{2em}引理}
\newtheorem{Corollary}[Theorem]{\hspace{2em}推论}
\newtheorem{Proposition}[Theorem]{\hspace{2em}命题}
\newtheorem{Definition}[Theorem]{\hspace{2em}定义}
\newtheorem{Example}[Theorem]{\hspace{2em}例}

% 列表格式
\usepackage{enumitem}
\setlist{nosep}

% 编号格式
\usepackage{chngcntr}
\counterwithin{table}{section}
\counterwithin{figure}{section}
\counterwithin{equation}{section}

% 页面风格
\usepackage{fancyhdr}
\pagestyle{fancy}

% 其他
\usepackage{booktabs} % 三线表
\usepackage{graphicx} % 图片导入
\hypersetup{ % 超链接样式
  colorlinks=true,    
  pdfborder=0 0 1,
  citecolor=black,
  linkcolor=black,
}
\usepackage{hologo} % logo
\usepackage{xurl}

% 引用编号格式
\labelformat{section}{第\zhnumber{#1}章}
\labelformat{subsection}{第\zhnumber{#1}节}
\labelformat{figure}{图~#1\space}

% 文档元数据
\title{南开大学本科生毕业论文模板\\ v3.0}
\title*{Released in May, 2023}
\studentid{1234567}
\author{XXX}
\grade{2019 级}
\major{计算机科学与技术}
\department{计算机科学与技术}
\college{计算机学院}
\adviser{XXX\quad 教授 \\ 罗翔\quad 教授 \\（张三大学）}
\date{\today}
\keywords{南开大学；毕业论文；模板}
\keywords*{dark energy model; dynamic analysis}

\begin{document}

\pagenumbering{Roman}

\maketitle
\begin{titlepage}
  \declaration
\end{titlepage}

\begin{abstract}
此模板最初按照南开大学计算机/网络空间安全学院本科生毕业设计要求制作整理。我们希望此模板能够帮助更多的同学。本模板会根据同学们的需求进行更新。最新版本请关注本模版主页：\href{https://tr0py.github.io/NKU-thesis-template-2020/}{NKThesis: 南开大学本科生毕业论文模板}。

在2023年5月，本模板经历了一次较大的重构，根据\href{http://jwc.nankai.edu.cn/2022/1124/c24a497818/page.htm}{2022年11月版《南开大学本科毕业论文（设计）指导手册》}重新进行了实现，原有接口不再兼容。
  
如果遇到使用问题，你可以用以下方式求助：
\begin{itemize}
  \item 在\href{https://github.com/Tr0py/NKU-thesis-template-2020/discussions}{GitHub讨论区}询问你不懂的地方。
  \item 在\href{https://github.com/Tr0py/NKU-thesis-template-2020/issues}{GitHub Issue}中提出你遇到的待修复的问题。
  \item 如果不会使用GitHub，你可以在QQ群中提问：群号469868290。
\end{itemize}

感谢历届的贡献者们：
\begin{itemize}
    \item 安祺，23届本科毕业生，\url{aqnin@outlook.com}
    \item 费迪，22届本科毕业生，\url{fd@nbjl.nankai.edu.cn}
    \item 林雪，22届本科毕业生，\url{ylsnowy@gmail.com}
    \item 黎鸿儒，22届本科毕业生，\url{lihongru@mail.nankai.edu.cn}
    \item 潘宇，20届本科毕业生，\url{panyu@nbjl.nankai.edu.cn}，\href{https://nbjl.nankai.edu.cn/2019/0513/c19244a264683/page.htm}{主页}
    \item 赵子懿，20届本科毕业生，\url{troppingz@gmail.com}，\href{https://tr0py.github.io/}{主页}
\end{itemize}
\end{abstract}
  
\begin{abstract*}
Since 1998, two independent supernova research groups have discovered that the universe is accelerating, and the dark energy has become a hot topic in cosmology. 
\end{abstract*}

\tableofcontents

\pagenumbering{arabic} 

\section{介绍}

\subsection{维护中}

请follow我们最新的进度 \url{https://github.com/Tr0py/NKU-thesis-template-2020}, 
以获取最新的模板bug修复及更新。如有问题请在GitHub留言，我们会尽快回复。
欢迎加入毕业论文模板维护的工作中。

\subsection{免责声明}

本模板本模板参考\href{http://jwc.nankai.edu.cn/2022/1124/c24a497818/page.htm}
{2022年11月版《南开大学本科毕业论文（设计）指导手册》}编写。

您自愿使用这个模板。
提供本模板的目的是为了给您的论文写作带来方便，然而，
作者不保证这个模板完全符合学校的要求，也不对由此产生的任何后果负责。
如果您不同意这些条款，请不要使用这个模板。

\section{使用说明} \label{chpt:A}

\subsection{引用章节号}
\label{sec:ex:A}

引用章节号请参考如下格式: \ref{chpt:A}\ref{sec:ex:A}

\subsection{中英文间隔}

使用 \hologo{XeLaTeX}\ 编译时，会自动在中英文转换时添加必要的空格。

\subsection{图表}

插图示例，如\ref{fig:nku}所示。

\begin{figure}
\centering
\includegraphics[viewport=0 0 2984 969,width=40mm]{figures/nankaidaxue.pdf}
\caption{南开大学}
\label{fig:nku}
\end{figure}

表格的例子：

\begin{table}
  \centering
  \caption{NKThesis 预调用的宏包}
  \begin{tabular}{l|l}
  \hline
  编译方式 & 调用的宏包\\ \hline
  \hologo{XeLaTeX} & xeCJK, CJKnumb, graphicx, mathptmx \\ \hline
  [PDF]\LaTeX & CJK, CJKnumb, uniGBK, graphicx, mathptmx \\
  \hline
  \end{tabular}
  \end{table}

\subsection{字体}

一般情况下, 您不需要显式地设置字体. 如果确实需要, 请使用以下命令

\begin{verbatim}
宋体:  \rmfamily\upshape 或 \songti
黑体:  \bfseries 或 \heiti
楷体:  \itshape  或 \kaiti
仿宋:  \ttfamily 或 \fangsong
加粗:  \jiacu
\end{verbatim}

\subsection{参考文献} \label{manual:ref}
参考文献引用:
\cite{ChenCheChen2001,Nadkarni-1992,Hua-Wang-1973}
文献\cite{ZhuKeZhen,Huo}\cite{JiangXiZhou,Timoshenko,Zhang-Wang}
\cite{ZhuKeZhen}
测试\parencite{PQEMU}或者\parencite{li1988ivy}

第一次编译\cite{hong2012hqemu}\cite{zaharia2010spark}\cite{wenwenasplos2018}

\subsubsection{录入参考文献}

本模板采用 biblatex 宏包管理参考文献。如果你对此不熟悉，可以参考biblatex使用手册。

\subsection{一些建议}
\subsubsection{关于分数的写法}

\LaTeX 提供宏命令\verb+\frac+, 用以打印分数. 为使得版面整齐, 该命令的使用应遵循以下原则:

\begin{enumerate}
\item 仅在分行表达式中使用,
\item 不嵌套使用,
\item 不在上下标中使用.
\end{enumerate}

也就是说, 行内表达式和上下标中出现分数时一律用 $a/b$表示, 如
$(x+2)/((3x^2+4)(7+y))$. 下面是居中表达式:

\[
 x^2 = y^{1/2} +3.
\]

多行表达式: 尽量在加、减、乘、等号前换行. 在乘号前换行时,
下一行首用 \string\times:
\def\iint{\mathop{\int\!\!\!\int}}\def\calG{\mathcal G}
\begin{eqnarray}
&&\left|(W_{\psi_1}f)(a,b)-(W_{\psi_1}f)(a_j,b_{j,k})\right|^{2}\nonumber\\
&=&\frac{1}{C^{2}_{\varphi}}\Bigg|\iint_{\calG} (W_{\varphi}f)(s,t) \nonumber\\
&&\qquad\times \Bigg( (W_{\psi_1}\varphi)\left(\frac{a}{s},
\frac{b-t}{s}\right)
     -(W_{\psi_1}\varphi)\left(\frac{a_{j}}{s}, \frac{b_{j,k}-t}{s}\right)\Bigg)
  \frac{dsdt}{s^{d+1}}\Bigg|^2 \nonumber\\
&\le& \frac{1}{C^2_{\varphi}} \iint_{\calG} |(W_{\varphi}f)(s,t)|^2 \nonumber\\
&&\qquad \times\left| (W_{\psi_1}\varphi)\left(\frac{a}{s},
\frac{b-t}{s}\right)
    -(W_{\psi_1}\varphi)\left(\frac{a_{j}}{s}, \frac{b_{j,k}-t}{s}\right)\right|
   \frac{dsdt}{s^{d+1}}  \nonumber\\
&&\qquad \times   \iint_{\calG}\!
 \left|(W_{\psi_1}\varphi)\left(\frac{a}{s}, \frac{b-t}{s}\right)
    -(W_{\psi_1}\varphi)\left(\frac{a_{j}}{s}, \frac{b_{j,k}-t}{s}\right)\right|
 \frac{ ds dt}{s^{d+1}} \nonumber\\
&=& \frac{1}{C^2_{\varphi}} ....  \label{eq:a0}
\end{eqnarray}


\subsubsection{标点}
科技文献中一般用半角标点, 请参考《中国科学》发表的论文.

如果使用全角标点, 可以使用
\begin{verbatim}
  \punctstyle{<style>}
\end{verbatim}
选择标点样式, 有效值为
\begin{verbatim}
  quanjiao (所有标点符号占一个汉字宽度，
            相邻两个标点占一个半汉字宽度)
  banjiao  (所有标点符号占半个汉字宽度)
  hangmobanjiao (所有标点符号占一个汉字宽度，行末行首半角)
  kaiming  (句号、叹号、问号占一个汉字宽度，其他标点占半个汉字宽度)
\end{verbatim}
缺省为全角式。注意：不论选择哪种样式，都提供行末对齐(margin kerning)功能。



\begin{Theorem} \label{thm:latex}
\LaTeX 的输出是最完美的.
\end{Theorem}

先证明一个引理
\begin{Lemma} \label{thm:tex}
\TeX 文件在不同操作系统下的排版结果完全一致.
\end{Lemma}

\newcommand{\proof}[1][证明]{\par\textbf{#1.}\,\, }
\def\endproof{\hfill{\usefont{U}{msa}{m}{n}\char"03}}

\begin{proof}
这是证明.
\end{proof}


\begin{proof}[定理~\ref{thm:latex}的证明]
显然是错的.
\end{proof}

单个带编号的表达式
\begin{equation}\label{eq:a1}
x=y+z
\end{equation}

单个不带编号的表达式
\[
y=x-z.
\]

不带编号的多行表达式
\begin{eqnarray*}
x&=&y+z \\
 &=&z-s\\
 &<& 3. \\
 && \mbox{一些注释}
\end{eqnarray*}

带编号的多行表达式
\begin{eqnarray}
 x&=& y-z, \label{eq:aa1}\\
 y&=& x+z, \nonumber \\
 z&=&y-x. \label{eq:aa2}
\end{eqnarray}

引用:   定理\ref{thm:latex}的推论是什么呢？
方程式编号:  由(\ref{eq:a1})(\ref{eq:aa2})式.

\subsubsection{列举环境:  enumerate}

TODO 将会重写

\appendix
单标题附录\footnote{脚注}

\appendices
\section{标题}

\begin{reference}
\printbibliography[heading=none]
\end{reference}

\begin{acknowledgement}
  感谢使用本模板。
\end{acknowledgement}

\end{document}
